% Generated draft: requires mathematical and source review.
\documentclass[11pt]{article}
\usepackage[margin=25mm]{geometry}
\usepackage{fontspec}
\setmainfont{Latin Modern Roman}
\usepackage{amsmath,amssymb,mathtools,unicode-math}
\setmathfont{Latin Modern Math}
\usepackage{enumitem,needspace,xurl,hyperref}
\setlength{\parindent}{0pt}
\setlength{\parskip}{5pt}
\setlength{\emergencystretch}{3em}
\newcommand{\N}{\mathbb{N}}
\newcommand{\Z}{\mathbb{Z}}
\newcommand{\Q}{\mathbb{Q}}
\newcommand{\R}{\mathbb{R}}
\newcommand{\C}{\mathbb{C}}
\newcommand{\F}{\mathbb{F}}
\newcommand{\A}{\mathbb{A}}
\newcommand{\PP}{\mathbb{P}}
\newcommand{\E}{\mathbb{E}}
\newcommand{\eps}{\varepsilon}
\newcommand{\dd}{\,\mathrm{d}}
\begin{document}
\section*{UnsolvedMath ID: 30002477}
\textbf{Source problem code:} catalogue entry
\subsection{Definitions and mathematical statement}
Under the usual order-theoretic convention, a tree is a partially ordered set \(T=(X,\leq_T)\) such that, for each \(t\in X\), the strict predecessor set \(\{s\in X:s<_T t\}\) is well ordered by \(<_T\). A branch is a maximal chain. Write
\[
\mathcal T_{\aleph_1}=\{T:|T|=\aleph_1\text{ and }T\text{ has no branch of cardinality }\aleph_1\}.
\]
Here the displayed class records the size and branch restrictions explicitly supplied in the statement; any additional intended tree conventions remain unresolved.

For a specified notion \(\preccurlyeq\) of tree embeddability, the corresponding universality number is
\[
u_{\preccurlyeq}(\mathcal T_{\aleph_1})=
\min\bigl\{|\mathcal U|:\mathcal U\subseteq\mathcal T_{\aleph_1},\ 
\forall T\in\mathcal T_{\aleph_1}\ \exists U\in\mathcal U\ (T\preccurlyeq U)\bigr\},
\]
when this minimum is defined. Trees may be represented on the underlying set \(\omega_1\) to make this a set-sized definition, provided the chosen embeddability notion is invariant under isomorphism.

\textbf{Source question, with unresolved parameters explicit.} Is there a model of the intended background set theory in which
\[
\exists U\in\mathcal T_{\aleph_1}\ \forall T\in\mathcal T_{\aleph_1}\quad T\preccurlyeq U,
\]
that is, in which \(u_{\preccurlyeq}(\mathcal T_{\aleph_1})=1\)? Cardinalities, trees, and quantifiers in this assertion are interpreted internally to the model.

The source does not specify the embedding relation or the background axioms. It also does not explicitly require height \(\omega_1\) or countable levels. The supplied title refers to Aronszajn trees, whose usual definition includes those further requirements, so the intended class cannot be recovered from the title alone. Consequently, the displayed question is a parameterized reconstruction requiring clarification, rather than a uniquely determined formal proposition.
\subsection{Short English statement}
The question asks whether some model of set theory contains a single tree universal for all trees of cardinality \(\aleph_1\) with no branch of that cardinality. The embedding notion, background axioms, and intended tree class remain unspecified; in particular, the title suggests additional Aronszajn-tree restrictions absent from the statement.
\subsection{Sources}
\begin{itemize}
\item Local UnsolvedMath registry, supplied complete source record, permanent ID \(30002477\), exact source problem code \texttt{catalogue entry}; title: \emph{Universal Aronszajn Trees at Aleph One}. The record belongs to the Oberwolfach Reports Open Problems collection and gives proposed year 2014; no individual proposer is supplied.
\item Supplied report citation: \emph{Set Theory} (2014), Oberwolfach Reports, pp.~91--144, DOI \texttt{10.4171/owr/2014/02}. \url{https://doi.org/10.4171/owr/2014/02}. The supplied background describes Oberwolfach Reports as documenting mathematical workshops and their open-problem discussions. No exact problem location within the report is supplied.
\item Supplied literature source, labeled ``Oberwolfach set theory report'': \url{https://ems.press/content/serial-article-files/46488}.
\item Literature assessment dated \texttt{2026-08-21}: supplied dated curation, not independently verified. It labels the problem open under a triage rule because the extracted record was assessed as not being a well-formed identifiable set-theoretic proposition and no recoverable statement or resolution was available. This label does not establish the mathematical open status of any particular completed formulation.
\item The record separately labels its statement \texttt{corrected\_verified} and reports a manual semantic, typographic, and LaTeX review claiming a self-contained, meaning-preserving cleanup with no unsupported source-dependent additions. This is supplied verification metadata, not independent verification of the formulation or its relationship to the cited report.
\end{itemize}
\end{document}
